\section{Real-data applications}\label{sec:realdata}
\subsection{Sleep-EDF recordings and analysis variables}
We analyze the two studies in the expanded Sleep-EDF database, available from PhysioNet at \url{https://physionet.org/content/sleep-edfx/1.0.0/}. The Sleep Cassette study contains 153 recordings from 78 people, with two recordings for most participants, and was designed to study age-related sleep differences. The Sleep Telemetry study contains paired temazepam and placebo recordings from 22 people with mild difficulty initiating sleep \citep{Mourtazaev1995,Kemp2000}. The cassette recordings span extended day--night periods, and the stage-defined sampling below uses the whole available record. The age analysis asks whether temporal dependence at different amplitude quantiles varies with age within a sleep stage. The paired analysis asks how this dependence differs between medication and placebo nights for the same person. Ordinary spectral regression provides a parallel assessment of amplitude-sensitive power, so the two representations address complementary aspects of sleep EEG.

The bipolar EEG channels Fpz--Cz (frontopolar--central) and Pz--Oz (parietal--occipital) are sampled at 100 Hz. We use the supplied expert sleep-stage annotations and analyze stage 2 and rapid eye movement (REM) sleep separately. Eligible segments are nonoverlapping 60-second intervals contained wholly within one scored stage. Up to 20 segments per stage and night are selected at equally spaced positions in the ordered set of eligible intervals. Selection uses the annotations and recording times. For each channel, segments are excluded if they contain nonfinite values, have a standard deviation below $0.5\,\mu\mathrm V$, have a maximum absolute deviation from their median above $500\,\mu\mathrm V$, or have at least 1\% of samples at a digital recording limit.

The primary analysis for each study uses stage-2 Fpz--Cz recordings. REM and Pz--Oz analyses assess variation by sleep stage and recording location. The quantile grid is $\{0.1,0.2,\ldots,0.9\}$ and the frequency grid is 0.5--30 Hz in increments of 0.5 Hz. A Gaussian bandwidth of 1.5 Hz and $\epsilon=10^{-6}$ are used throughout the main analyses. Ordinary log-power regression uses the same segments, covariates, frequency grid, and smoothing bandwidth. Sensitivity analyses use bandwidths of 1 and 2 Hz and the first 30 seconds of each selected interval. All segments and nights from the same person remain in the same cluster.

For the age analysis, covariates are age in decades centered at 60 years, a female indicator, an indicator for the second recording night, elapsed time in hours since the first scored sleep epoch centered at three hours, and the square of this centered time. Each person has total weight one, with equal weight assigned to their available nights. The age coefficient is a linear adjusted association over the study's age range. For the paired medication analysis, spectral summaries are averaged within each night and the placebo-night mean is subtracted from the temazepam-night mean. These differences are regressed on an intercept, the difference in night number, and the difference in mean elapsed time of the selected segments. The intercept describes the adjusted paired contrast at equal elapsed time and a balanced recording order.

Tests and confidence bands use 4,999 bootstrap draws for the main analyses and 1,999 for the resolution sensitivity analyses. Confidence bands cover the specified frequency--quantile grid within each displayed surface. Global $p$-values for the secondary stage and channel analyses are reported as descriptive evidence for those separate questions. We also test the direct differences between nonmedian quantile coefficient curves and the median curve, using a maximum over all eight nonmedian quantiles and all frequencies. For these heterogeneity tests, Holm adjustment is applied across the four stage--channel combinations within each study. This distinguishes evidence of a general spectral association from evidence that its size varies with quantile level.
